%titolo{In fondo al mar (La sirenetta)}
%autore{Menken, Ashman}
%album{}
%tonalita{LA}
%famiglia{Altre}
%gruppo{}
%momenti{}
%identificatore{in_fondo_al_mar}
%data_revisione{2026_10_20}
%trascrittore{Francesco Endrici}
%video{}
\beginsong{In fondo al mar (La sirenetta)}[by={Menken, Ashman}]

\beginverse
\[A] Le alghe del \[E]tuo vi\[A]cino ti sembran più \[E]verdi \[A]sai,
vorresti andar \[E]sulla \[A]terra, non sai che gran \[E]sbaglio \[A]fai.
Se \[D]poi ti guardassi in\[A]torno, ve\[E]dresti che il nostro \[A]mar
è \[D]pieno di mera\[A]viglie, che \[E]altro tu vuoi di \[A]più?
\endverse

\beginchorus
In fondo al \[D]mar, \[A] in fondo al \[E]mar, \[A]
tutto ba\[D]gnato è molto \[E]meglio, credi a \[A]me. \[A7]
Quelli las\[D]sù che sgobba\[E]no 
sotto a quel \[F#-]sole svengo\[B]no 
mentre col \[D]nuoto ce la spas\[E]siamo in fondo al \[A]mar.
\endchorus

\beginverse
\chordsoff
Quaggiù tutti sono allegri guizzando di qua e di là,
invece là sulla terra il pesce è triste assai:
rinchiuso in una boccia, che brutto destino avrà:
se all'uomo verrà un po' fame, il pesce si papperà.
Oh, no!
\endverse

\beginchorus
In fondo al \[D]mar, \[A] in fondo al \[E]mar, \[A]
nessuno ci \[D]frigge o ci cu\[E]cina in fricas\[A]sea \[A7]
e non si \[D]rischia di abboc\[E]car, no, non c'è un \[F#-]amo in fondo al \[B]mar,
la vita è \[D]ricca di bolli\[E]cine
in fondo al \[D]mar, \[A] (in fondo al mar).
In fondo al \[E]mar, \[A] (in fondo al mar).
Con questo \[D]ritmo la vita è \[E]sempre dolce co\[A]sì \[A7]
(sempre così).
Anche la \[D]razza ed il sal\[E]mon,
sanno suo\[F#-]nare con pas\[B]sion,
qui c'è la \[D]grinta e ogni con\[E]certo è un succes\[A]son.
\endchorus

\beginverse
Il \[E]sarago suona il \[A]flauto, la carpa \[E]l'arpa, la platessa il \[A]basso,
poi c'è la \[D]tromba del pesce \[A]rombo, voilà, il \[E]luccio è il re nel \[A]blues.
La razza con il nasello, al violoncello con la sardina
all'oca\[D]rina e con l'o\[A]rata vedrai che \[E]coro si fa\[A]rà!
\endverse

\beginchorus
Sì, in fondo al \[D]mar \[A] (in fondo al mar),
in fondo al \[E]mar \[A] (in fondo al mar)
se la sar\[D]dina fa una mo\[E]ina c'è da impaz\[A]zir \[A7] (c'è da impazzir!).
Che c'è di \[D]bello poi las\[E]sù?
La nostra \[F#-]banda vale di \[B]più,
ogni mol\[D]lusco sa improvvi\[E]sare in fondo al \[A]mar,
ogni lu\[D]maca si fa un bal\[E]letto in fondo al \[A]mar,
e tutti \[D]quanti ci diver\[E]tiamo
qui sotto l'\[D]acqua, in mezzo al \[E]fango,
ah, che for\[D]tuna vivere in\[E]sieme in fondo al \[A]mar!
\endchorus
\endsong

